\documentclass{article}
\usepackage[preprint,nonatbib]{neurips_2025}
\usepackage{amsmath,amssymb,amsthm,booktabs,graphicx,hyperref,enumitem,longtable,array}
\usepackage{microtype}
\usepackage{setspace}
\setstretch{1.0}

\title{KernelWeave: Self-Compiling Skill Kernels for Language Models}
\author{Anonymous}

\newtheorem{theorem}{Theorem}
\newtheorem{definition}{Definition}
\newtheorem{proposition}{Proposition}

\begin{document}
\maketitle

\begin{abstract}
Large language models repeatedly solve the same families of tasks, yet most systems discard that success after each interaction. We introduce KernelWeave, a self-compiling architecture that converts verified task traces into reusable skill kernels. A kernel is not a cached answer or a prompt template; it is an executable object containing schemas, preconditions, postconditions, evidence requirements, rollback rules, and regression tests. At inference time, the runtime searches for a matching kernel, executes it when the evidence supports reuse, and falls back to raw generation when the match is weak. The result is a model system that accumulates durable competence instead of starting over on every prompt. We formalise the induction objective, derive a fail-closed reuse rule, present a working prototype, and define an evaluation programme that measures reuse, safety, and compounding capability.
\end{abstract}

\section{Introduction}

Current LLM systems improve through retrieval, routing, tool use, and fine-tuning, but they still have a basic flaw: when a task is solved well, the system often forgets how it solved it. The next similar request is answered again from scratch, paying the same reasoning cost and redoing the same safety checks. That is wasteful, and it is the real bottleneck for many practical deployments.

KernelWeave is built around a different idea. When a task family is solved successfully enough, the system should compile that success into a reusable computational object. On future prompts, the system should try the compiled skill first, rather than regenerating the entire solution path. In other words, the model should not merely remember text; it should accumulate verified competence.

The contribution is not a wrapper around a model. It is a compilation architecture for competence itself.

\paragraph{Contributions}
\begin{enumerate}[leftmargin=1.2em]
\item We define a skill kernel as a typed, executable, evidence-carrying representation of a solved task family.
\item We introduce a kernel induction objective that scores usefulness, evidence quality, reuse value, and failure risk.
\item We derive a fail-closed runtime rule that refuses reuse when evidence or regression checks are insufficient.
\item We present a full prototype with a kernel compiler, kernel store, runtime selector, CLI, and regression-gated sample kernels.
\item We outline an evaluation plan that measures reuse rate, compounding savings, drift robustness, and safety.
\end{enumerate}

\section{Problem statement}

Let a prompt family be denoted by $\mathcal{P}$, and let successful traces over that family be $\tau \in \mathcal{T}(\mathcal{P})$. We seek a compilation operator
\[
\mathcal{K}: \mathcal{T}(\mathcal{P}) \to \mathbb{K},
\]
where the resulting object is executable, testable, and reusable.

A kernel must satisfy six requirements: it must be typed, it must define preconditions and postconditions, it must carry evidence requirements, it must fail closed when evidence is insufficient, it must be rollbackable, and it must be stored as a durable artifact.

\begin{definition}[Skill kernel]
A skill kernel is a structured, executable representation of a solved task family, containing input and output schemas, a step plan, evidence constraints, rollback logic, and regression tests.
\end{definition}

This makes the unit of reuse concrete. A kernel is not a note. It is a validated capability module.

\section{KernelWeave framework}

We represent a task family by prompts $p \in \mathcal{P}$ and a successful trace by
\[
\tau = (e_1, e_2, \dots, e_n),
\]
where each event $e_i$ is one of: plan, tool call, observation, evidence item, verification, decision, or failure.

A kernel $k$ is a structured object
\[
 k = (I_k, O_k, S_k, R_k, E_k, X_k, \pi_k, c_k),
\]
where $I_k$ is the input schema, $O_k$ the output schema, $S_k$ the step plan, $R_k$ the rollback policy, $E_k$ the evidence requirement set, $X_k$ the regression suite, $\pi_k$ the execution policy, and $c_k$ the confidence score.

The important part is that a kernel can be executed, checked, and revoked. It is therefore not a static memory record.

\subsection{Induction objective}
Given a trace $\tau$, the compiler seeks a kernel that maximises
\[
J(k;\tau)=\lambda_u U(k;\tau)+\lambda_e Q(k;\tau)+\lambda_r V(k)-\lambda_f F(k)-\lambda_o O(k),
\]
where $U$ measures task utility recovered from the trace, $Q$ measures evidence quality, $V$ measures expected reuse value, $F$ measures failure risk, and $O$ measures operational overhead.

A useful concrete form is
\[
U(k;\tau)=\alpha A(k,\tau)+\beta E(k)+\gamma C(k),
\]
where $A$ is applicability, $E$ is evidence completeness, and $C$ is compaction of the trace into executable structure.

\subsection{Runtime score}
At inference time, a prompt $p$ is matched against kernels by
\[
S(p,k)=w_m M(p,k)+w_c c_k+w_e \mathrm{cov}(k,p)-w_f \mathrm{risk}(k),
\]
where $M$ is semantic match, $\mathrm{cov}$ is evidence coverage for the prompt, and $\mathrm{risk}$ captures drift and overfitting. The runtime executes the kernel only if $S(p,k) \ge \theta$ and the regression gates pass.

\section{Algorithm}

KernelWeave compiles behaviour in five steps.

\begin{enumerate}[leftmargin=1.2em]
\item Observe a successful trace for a task family.
\item Extract events, tools, evidence, and output constraints.
\item Infer a typed schema and a step plan.
\item Attach rollback rules, evidence requirements, and tests.
\item Admit the kernel only if its confidence and regressions meet threshold.
\end{enumerate}

\begin{proposition}[Compile-or-reject]
A candidate kernel is either admitted with a recorded confidence and test suite, or rejected entirely. There is no partial admission.
\end{proposition}

This matters because half-baked kernels are worse than no kernel. They create the illusion of reuse while quietly increasing failure probability.

\subsection{Fail-closed reuse}
\begin{definition}[Evidence debt]
Evidence debt is the gap between the evidence required by a kernel and the evidence currently available for the active prompt.
\end{definition}

If evidence debt is non-zero, the runtime must either request more evidence or fall back to generation.

\begin{theorem}[Fail-closed reuse]
If a kernel fails its regression suite or its evidence requirements are not satisfied, then the runtime does not execute it.
\end{theorem}

\begin{proof}
The runtime only admits kernels whose stored confidence exceeds threshold and whose regression tests pass against the current constraints. If evidence is missing, the evidence debt is positive, so the gating condition fails. If the regression suite fails, the kernel is flagged invalid. In either case the execution branch is not taken, and the system falls back to raw generation or a more general kernel. Hence reuse is fail-closed.
\end{proof}

\subsection{Compounding benefit}
If a task family occurs with probability $\rho$ and solving it from scratch costs $C_s$ compute units, while executing the kernel costs $C_k$ with induction overhead $C_i$, then expected net savings after $T$ occurrences is approximately
\[
\Delta(T)=T\rho(C_s-C_k)-C_i.
\]
The kernel is worthwhile when $\Delta(T)>0$. That is the basic compounding argument: once recurrent tasks are frequent enough, compilation amortises.

\section{Prototype implementation}

We built a working prototype rather than a mockup.

The system includes a kernel store persisted on disk, a compiler that turns traces into kernels, a runtime that scores prompts against the kernel store, a CLI for initializing and planning, JSON-backed traces and kernels, sample kernels with deterministic scoring, and regression gating.

The prototype already supports the core cycle: trace $\rightarrow$ kernel $\rightarrow$ selection $\rightarrow$ fallback.

\subsection{Prototype kernels}
The sample store includes kernels for task families such as safe shell command synthesis and artifact comparison. Those are intentionally different: one is safety-heavy and one is comparison-heavy. The point is to show that the compiler can induce kernels from diverse task structures, not just one narrow benchmark.

\subsection{Implementation sketch}
The kernel compiler scores the trace for evidence density and tool diversity, generates a kernel ID from task family and summary, persists the kernel, and later scores prompts against all kernels. If the best match exceeds threshold, the runtime executes the kernel; otherwise it falls back to generation.

\section{Evaluation design}

A serious evaluation should measure reuse, not just answer quality.

\begin{enumerate}[leftmargin=1.2em]
\item \textbf{Reuse rate}: how often the runtime chooses a kernel instead of raw generation.
\item \textbf{Kernel accuracy}: when a kernel is selected, what fraction of runs satisfy the output contract and the evidence contract.
\item \textbf{Compounding savings}: how much compute is saved over time compared with always generating from scratch.
\item \textbf{Drift robustness}: how often slightly shifted wording still maps to the correct kernel.
\item \textbf{Safety calibration}: how often the runtime correctly rejects stale or under-evidenced kernels.
\end{enumerate}

A useful aggregate metric is
\[
\chi = \frac{\text{saved cost over horizon} - \text{induction cost}}{\text{induction cost}},
\]
which becomes positive when the system saves more than it spent to compile the kernels.

\section{Discussion}

KernelWeave is deliberately more ambitious than a memory wrapper and less magical than a general intelligence claim. The thesis is that LLMs can become better systems if solved behaviour is compiled into durable, executable structure.

This matters in several ways. Repeated enterprise tasks can become cheaper and more consistent. Safety-critical flows can accumulate explicit rollback logic. Benchmarks can measure compounding capability rather than one-off performance. Models can become more specialised without retraining from scratch every time.

The idea is not that every prompt becomes a kernel. The idea is that the system learns where the stable structure is and preserves it.

\section{Limitations}

There are real failure modes.
\begin{itemize}[leftmargin=1.2em]
\item Prompt families can be fuzzy and kernel boundaries can blur.
\item Kernels can overfit traces if evidence is too narrow.
\item Similarity scoring can miss semantically related prompts.
\item Behaviour that is genuinely creative or open-ended may not compress into kernels well.
\item Evidence can become stale as external tools, policies, or data change.
\end{itemize}

These are not reasons to dismiss the architecture. They are the actual research questions it creates.

\section{Conclusion}

KernelWeave reframes an LLM as a system that can compile its own successful behaviour into reusable skill kernels. That is the core novelty: competence becomes a persistent object with schemas, tests, rollback, and confidence, rather than a transient answer hidden in a chat log. If the evaluation confirms the reuse and compounding properties, the result would be a different class of LLM system: one that gets better not only by scaling, but by accumulating verified structure.

\appendix
\section{Compilation pseudocode}

\begin{enumerate}[leftmargin=1.2em]
\item Input: successful trace $\tau$, task family label $f$, description $d$.
\item Score evidence density, tool diversity, and trace compression value.
\item Construct $(I_k,O_k,S_k,R_k,E_k,X_k,\pi_k,c_k)$.
\item Store the kernel and its source trace ID.
\item At runtime, search kernels by prompt similarity.
\item If the best kernel exceeds the threshold and passes gates, execute it; otherwise fall back.
\end{enumerate}

\section{Additional derivation}
If the expected occurrence count for a task family is $N$ and the average from-scratch cost is $C_s$, then the expected total cost without kernels is $N C_s$. With kernels it is approximately $N C_k + C_i$, where $C_k < C_s$ if the kernel is useful. Therefore the break-even point is
\[
N > \frac{C_i}{C_s-C_k}.
\]
This is the cleanest reason the approach can scale: once task families repeat, compilation amortises.

\end{document}
